%! Linear Functions: Slope, Forms & Graphing — Unit 1, Lesson 1.1 — Homework (Answer Key)
% Mirrors homework/main.tex; each \writelines{2} is answered with exactly two \ansline's.
\documentclass[10pt]{article}
\usepackage{algebra2-article}
\usepackage{algebra2-key}   % pulls in -boxes; do NOT also load -boxes
% Coordinate grid: {xmin}{xmax}{ymin}{ymax}
\newcommand{\lgrid}[4]{%
  \draw[step=1, linegray!40, very thin] (#1,#3) grid (#2,#4);
  \draw[->, semithick, charcoal] (#1,0)--(#2,0) node[below right, font=\scriptsize]{$x$};
  \draw[->, semithick, charcoal] (0,#3)--(0,#4) node[left, font=\scriptsize]{$y$};}

\begin{document}

\pageheader{Unit 1, Lesson 1.1}{Homework --- Answer Key}

\vspace{0.10in}

\begin{notesbox}{Practice}
Show your work. Remember: \emph{slope-intercept} hands you a rate and a start, \emph{point-slope} is
built from a slope and \emph{any} point, and \emph{standard} is the tidy integer version --- all
three can describe the same line.

\begin{enumerate}[leftmargin=1.9em, itemsep=11pt]
  \item \textbf{Slope from two points.} \\[3pt]
        Through $(-2,7)$ and $(4,-5)$: $m = \ans{$-2$}$ \quad
        Through $(-3,4)$ and $(6,4)$: $m = \ans{$0$}$ \quad
        Through $(5,1)$ and $(5,-3)$: $m = $ \ans{undefined}

  \boxguard[18]
  \item \textbf{One line, two ways.} The line shown passes through $(0,-1)$ and $(3,5)$.
    \begin{center}
    \begin{minipage}[c]{0.32\linewidth}\centering
    \begin{tikzpicture}[scale=0.30]
      \lgrid{-3}{5}{-6}{7}
      \draw[very thick, forest] (-2.5,-6)--(4,7);
      \fill[charcoal] (0,-1) circle (4pt);
      \fill[charcoal] (3,5) circle (4pt);
    \end{tikzpicture}
    \end{minipage}\hfill
    \begin{minipage}[c]{0.64\linewidth}
    \begin{enumerate}[label=(\alph*), leftmargin=1.9em, itemsep=4pt]
      \item Slope: \ans{$2$} \; Slope-intercept form: \ans{$y=2x-1$}
      \item Point-slope form, using $(3,5)$: \ans{$y-5=2(x-3)$}
      \item Are (a) and (b) the same line? \ans{yes} \; Show it:
    \end{enumerate}
    \end{minipage}
    \end{center}
    \begin{work}
      y - 5 &= 2(x - 3) \\
      y - 5 &= 2x - 6 \\
          y &= 2x - 1
    \end{work}

  \item \textbf{Change its outfit.} Rewrite $y + 2 = -\tfrac34(x - 4)$ in slope-intercept form,
        then in standard form.
    \begin{work}
      y + 2 &= -\tfrac34(x-4) \\
      y + 2 &= -\tfrac34 x + 3 \\
          y &= -\tfrac34 x + 1 \\
         4y &= -3x + 4 \\
    3x + 4y &= 4
    \end{work}
    Slope-intercept: \ans{$y=-\tfrac34x+1$} \quad Standard: \ans{$3x+4y=4$}

  \boxguard[18]
  \item \textbf{The parent, and the two odd ones.} Start from $y = x$ and end at $y = -3x + 2$.
    \begin{enumerate}[label=(\alph*), leftmargin=1.9em, itemsep=4pt]
      \item Steeper or flatter? \ans{steeper} \; Flipped? \ans{yes} \;
            Shifted? \ans{up $2$}
      \item Slope of $y = 4$: \ans{$0$} \quad Slope of $x = 4$: \ans{undefined}
      \item Which of those two \emph{cannot} be written as $y = mx + b$, and why?
        \ansline{$x=4$. Its slope is undefined, so there is no number to put in for $m$;}
        \ansline{and every point on it has $x=4$, which $y=mx+b$ can never force.}
    \end{enumerate}

  \boxguard[20]
  \item \textbf{Model it.} A candle burns at a steady rate. After $1$ hour it is $17$~cm tall;
        after $4$ hours it is $11$~cm tall.
    \begin{enumerate}[label=(\alph*), leftmargin=1.9em, itemsep=4pt]
      \item Find the rate, then the rule.
        \begin{work}
          m &= \frac{11 - 17}{4 - 1} \\
            &= \frac{-6}{3} \\
            &= -2 \\
          h(t) &= -2t + 19
        \end{work}
      \item The $y$-intercept means \ans{the candle started at 19 cm} \; After $6$ hours:
            \ans{$7$ cm} \; Burns out at \ans{$t=9.5$ h}
    \end{enumerate}

  \item \textbf{Multiple choice (SOL-style).} Which line is \emph{perpendicular} to
        $y = \tfrac12 x - 4$ and passes through $(2,-1)$?
    \begin{enumerate}[label=(\Alph*), leftmargin=2.0em, itemsep=1pt]
      \item $y = \tfrac12 x - 2$
      \item \ans{$y = -2x + 3$}
      \item $y = 2x - 5$
      \item $y = -\tfrac12 x$
    \end{enumerate}
    Say in one sentence how you ruled out one of the others.
    \ansline{\textbf{B} is correct: $m=-2$ is the opposite reciprocal of $\tfrac12$, and $-2(2)+3=-1$.}
    \ansline{\textbf{A} has the \emph{same} slope --- that is parallel, not perpendicular.}
\end{enumerate}
\end{notesbox}

\boxguard[18]
\begin{extensionbox}
\textbf{Build and reason.}
\begin{enumerate}[label=(\alph*), leftmargin=1.9em, itemsep=6pt]
  \item Write the equation of the line through $(-4,3)$ that is \emph{parallel} to $2x - y = 5$.
        \begin{work}
          2x - y &= 5 \\
               y &= 2x - 5 \\
           y - 3 &= 2(x + 4) \\
               y &= 2x + 11
        \end{work}
  \item In class you wrote a line from one of its two points; a classmate started from the other and
        got an equation that looked different. Explain why point-slope gives the \emph{same} line no
        matter which point on the line you start from.
        \ansline{A line is fixed by its slope and \emph{one} of its points, and every point on it}
        \ansline{satisfies the same rule --- so expanding any of them gives one $y=mx+b$.}
\end{enumerate}
\end{extensionbox}

\begin{spiralbox}
\textbf{Coming next --- Lesson 1.2: Absolute Value Equations \& Inequalities.} Every line you wrote
today crosses zero exactly once. Next comes a rule that reaches the same value \emph{twice}, because
it measures \emph{distance} --- so one equation will have two answers.
\end{spiralbox}

\end{document}
